% generated by code/p5_calc/release_p19.py -- do not edit
\begin{center}\scriptsize\begin{tabular}{p{3.2cm}p{2.6cm}p{2.6cm}p{2.6cm}}\hline
assignment $(c,h)$ & defect pole & coeff. of L & local finite \\\hline
\texttt{Q4S (declared, R2)} & \texttt{0} & \texttt{0} & \texttt{0} \\
\texttt{QdS (hat)} & \texttt{0} & \texttt{0} & \texttt{0} \\
\texttt{c tilde, h hat} & \texttt{0} & \texttt{0} & \texttt{0} \\
\texttt{c tilde, h mixed} & \texttt{0} & \texttt{0} & \texttt{0} \\
\texttt{c hat, h mixed} & \texttt{0} & \texttt{0} & \texttt{0} \\
\texttt{all tilde} & \texttt{0} & \texttt{0} & \texttt{0} \\
\texttt{c hat, h tilde} & \texttt{0} & \texttt{0} & \texttt{0} \\
\hline\end{tabular}\end{center}
\noindent Route G over \texttt{Q(i)}: rank $\sigma_0|_D$ = 5, dim $D\cap\ker\sigma_0$ = 6, dim $D\cap\mathrm{Im}\,\sigma_0$ = 6, cohomology 0, quotient 5 (complement \texttt{h.h\#0, h.h\#1, h.h\#2, c.cb\#0, c.cb\#1}); rank $P_{ch}\sigma_0$ = 3 (Route C: 3).
Over Q with the re/im split (p14r convention): dim $D\cap\mathrm{Im}$ = 3, quotient 8, dim $D\cap\ker$ = 5.

\noindent Quotient coordinates of the $\Gamma_1$ local data: pole \texttt{['0', '0', '0', '-3/8', '1/4']}, coeff. of L \texttt{['0', '0', '0', '3/8', '-1/4']}, local \texttt{['0', '0', '0', '-2/3', '5/12']}.
\begin{center}\scriptsize\begin{tabular}{p{4.6cm}p{6.4cm}p{1.2cm}}\hline
mutation & outcome & changes \\\hline
\texttt{drop\_graph\_D03} & \texttt{OBSTRUCTION (nonlocal residual: the L-coeffi} & yes \\
\texttt{drop\_graph\_D01} & \texttt{OBSTRUCTION (nonlocal residual: the L-coeffi} & yes \\
\texttt{flip\_phase\_D03} & \texttt{OBSTRUCTION (nonlocal residual: the L-coeffi} & yes \\
\texttt{symmetry\_factor\_D03\_x2} & \texttt{OBSTRUCTION (nonlocal residual: the L-coeffi} & yes \\
\texttt{wrong\_projection\_trace\_operator\_hh} & \texttt{OBSTRUCTION (nonlocal residual: the L-coeffi} & yes \\
\texttt{wrong\_projection\_gauge\_fixed\_operator\_Kc} & \texttt{OBSTRUCTION (nonlocal residual: the L-coeffi} & yes \\
\texttt{naive\_4D\_D03\_evanescent\_dropped} & \texttt{OBSTRUCTION (nonlocal residual: the L-coeffi} & yes \\
\texttt{naive\_4D\_D13\_Gt\_dropped} & \texttt{NONZERO, CANCELLABLE by a counterterm in the} & yes \\
\hline\end{tabular}\end{center}
\noindent Counterterm solve: \texttt{ZERO}; \texttt{x = 0 (no counterterm needed beyond the MS-bar poles); the identity-preserving finite renormalisations are ker(P\_ch sigma\_0) cap D (dim 8), of which 2 modulo D cap Im sigma\_0}.

\noindent\textbf{Why zero.} \texttt{each surviving ch row vanishes separately in the projected slot; no counterterm cancellation and no cancellation between rows is involved: (i) hh row (D01+D02+D03+D11+D12): full-row transverse decoupling -- D01 and D03 are individually nonzero and cancel against each other in the slot (p17/p18 transverse sum), D11 is transverse by itself, D02 = D12 = 0; (ii) K.c row (D13): D13 reaches the ch slot only through T1 and E1 (T2, T3, E2 are longitudinal, E transverse) and A\_T1 = A\_E1 = 0 exactly in d; (iii) one-point row: the tadpoles D14a-c vanish under the declared prescription.  The defect is therefore zero without any counterterm (solve x = 0).}

\noindent\textbf{Flag definitions.} \textbf{\texttt{projected\_defect\_computed}}: \texttt{Delta = P\_ch[(S\_Psi, Gamma\_Psi)]\_\{hbar kappa\^{}2\} assembled from every evaluated input of the three surviving rows (none Missing), exact in d, pole + nonlocal + MS-bar local parts, all seven external assignments}. \textbf{\texttt{projected\_identity\_verified}}: \texttt{Delta = 0 for the REGULATED one-loop functional Gamma1: pole, nonlocal and local parts vanish for all seven assignments, exactly in d, with all controls true and every mutation changing the outcome}. \textbf{\texttt{projected\_restoration\_verified}}: \texttt{p20 definition (replaces p19's, which follows from the identity + locality): FULL sigma\_0 closure of the COMPLETE one-loop kappa\^{}2 MS-bar counterterm S\_ct = -(pole) in the declared domain, h.h + c.K + c.cb (ghost self-energy Sigma held fixed at its P15 value, orientation converted, derived relative factor +1), in every ghost +1 sector (h.c, c.b, c.cbs), together with projected\_identity\_verified.  Scope: the one-loop kappa\^{}2 counterterm only; NOT bulk symmetry, NOT O.1}. \textbf{\texttt{bulk\_symmetry\_restored}}: \texttt{hard stop: False regardless (not claimed; O.1 NOT\_CLOSED)}.

\noindent\textbf{K1 (corrected in p20; erratum to p19).} derived relative factor \texttt{1}; orientation ratio (route-G (c,cb) over P15 (cb,c)) \texttt{['-1', '-1']}; Sigma pole in (c,cb) coordinates \texttt{['3/8', '-1/4']}; full sigma\_0 image of the complete MS-bar counterterm: \texttt{\{'closed': True, 'entries': 0, 'sectors': []\}}. \texttt{ERRATUM to p19: the p19 statement 'closure fails with the derived factor, needs -1, unresolved' was a Grassmann-orientation bookkeeping error (P15 (cb,c) coordinates placed in route-G (c,cb) slots).  Corrected: derived +1, full closure holds (finding K1 of the p19 independent rebuild).}

\noindent\textbf{J1 (finding about the p14r contract).} \texttt{FINDING ABOUT THE p14r CONTRACT (contract NOT modified): the p14r counterterm-domain counts (Im sigma\_0 3, quotient 8, D cap ker 5, P\_ch sigma\_0 rank 4) are artefacts of a re/im split over Q (the densities carry explicit i for odd derivative counts, so i-multiples count as independent).  The counts over Q(i) -- sigma\_0 rank 5, D cap ker 6, D cap Im 6, cohomology 0, quotient 5, P\_ch sigma\_0 rank 3 (= Route C) -- equal the real dimensions under a compatible real structure.  The p14r certificate stays as released; both are reported}
